\section{Formalizzazioni dei concetti base}
\subsection{Algoritmo}
\subsubsection*{Definizione formale}
Un algoritmo è una macchina a stati finiti che riceve una sequenza di input, effettua una sequenza finita di transizioni tra
gli stati interni della macchina e produce una sequenza di output.

Si osserva che tale definizione coincide con la definizione di una macchina di Turing.

\subsubsection*{Stopping problem}
Si definisce lo ``stopping problem'' come la proprietà di un algoritmo di non poter continuare all'infinito. Ad esempio non
esiste un algoritmo che calcoli tutte le cifre di $\pi$ o che verifica se un'equazione diofantea (o diofantina) ha soluzioni
intere.

\subsubsection*{Progettazione di un algoritmo}
La progettazione di un algoritmo è il processo di ideazione e formalizzazione di un algoritmo per risolvere uno specifico
problema computazionale. Questo processo coinvolge:
\begin{itemize}
	\item la conoscenza del dominio del problema ed eventuali proprietà che lo caratterizzano, in quanto potrebbero
	semplificare il processo di risoluzione
	\item la conoscenza di metodi generali di progettazione di algoritmi, come i paradigmi di programmazione
\end{itemize}

\subsubsection*{Analisi di un algoritmo}
L'analisi di un algoritmo è il processo di valutazione dell'algoritmo secondo i seguenti aspetti:
\begin{itemize}
	\item correttezza: valuta se l'algoritmo porta alla soluzione corretta del problema per ogni istanza
	\item efficienza: valuta il consumo di risorse in termini di passi effettuati (complessità temporale), di memoria
	occupata (complessità spaziale) ed eventuali altri parametri (es. consumo energetico)
\end{itemize}

\subsection{Problema computazionale}
\subsubsection*{Definizione formale}
Un problema computazionale è composto dai seguenti tre elementi:
\begin{itemize}
	\item un insieme \(\mathcal{I}\) di istanze, che rappresentano i possibili input del problema, si nota che alcuni
	problemi potrebbero avere un numero infinito di istanze, per cui \(\mathcal{I}\) potrebbe essere un insieme infinito
	\item un insieme \(\mathcal{S}\) di soluzioni, che rappresentano gli output del problema per ogni istanza
	di \(\mathcal{I}\)
	\item una relazione \(\pi \subseteq \mathcal{I} \times \mathcal{S}\) che associa ogni istanza \(i \in \mathcal{I}\)
	ad una o più soluzioni \(s_j \in \mathcal{S}\) valide per quella istanza (è un sottoinsieme del prodotto cartesiano
	tra \(\mathcal{I}\) e \(\mathcal{S}\))
\end{itemize}

\subsubsection*{Rappresentazione delle istanze e delle soluzioni come stringhe}
I dati di un'istanza e di una soluzione sono rappresentati come stringhe, ovvero sequenze finite di simboli. Si assumono
le seguenti convenzioni:
\[\begin{array}{c c l }
	1. & \quad \Sigma = \{ 0, \; 1, \; \dots \} & \qquad \text{l'alfabeto dei simboli delle stringhe} \\[6pt]
	2. & \quad \Sigma^0 = \{\varepsilon\} & \qquad \text{l'insieme della sola stringa vuota} \\[6pt]
	3. & \quad \Sigma^i & \qquad \text{l'insieme di tutte le stringhe di lunghezza} \; i \; \text{(potenza cartesiana)} \\[6pt]
	4. & \quad \Sigma^+ = \bigcup_{i = 1}^{\infty} \Sigma^i & \qquad \text{l'insieme di tutte le stringhe esclusa la stringa vuota} \\[6pt]
	5. & \quad \Sigma^* = \bigcup_{i =0}^{\infty} \Sigma^i & \qquad \text{l'insieme di tutte le stringhe inclusa la stringa vuota} \\[6pt]
	6. & \quad \mathcal{I} \subseteq \Sigma^*, \quad \mathcal{S} \subseteq \Sigma^* & \qquad \text{gli insiemi di istanze e soluzioni sono sottoinsiemi di } \Sigma^*
\end{array}\]
La taglia di una istanza è la quantità di informazione necessaria per rappresentarla ed è pari al numero di simboli che
compongono la stringa che la rappresenta, o meglio la sua lunghezza. Si noti che per un numero intero \(n\) rappresentato
in binario, la taglia è il numero di cifre che lo compongono \(\approx \log_{2} n\) bit.

\subsubsection*{Osservazioni sulla rappresentazione dei dati nell'insieme delle istanze}
L'insieme delle istanze \(\mathcal{I}\) impone il tipo di rappresentazione con cui i dati in input devono essere codificati.
Fornire i dati in una diversa rappresentazione potrebbe portare a soluzioni errate, ad esempio se l'algoritmo svolge la somma
di numeri in base 10 e gli vengono forniti in base 2.

Inoltre la scelta della rappresentazione influisce sull'efficienza dell'algoritmo in quanto alcune operazioni risultano
più efficienti in una rappresentazione rispetto ad un'altra. Ad esempio la convoluzione nel dominio del tempo è \(O(n^2)\),
mentre nel dominio della frequenza è \(O(n \log n)\).

All'interno dell'algoritmo è possibile aggiungere una fase di conversione dei dati da una rappresentazione ad un'altra
(ad esempio per ridurre la complessità computazionale di alcuni problemi), ma è sempre necessario specificare la
rappresentazione dei dati in ingresso.

\subsubsection*{Funzioni e algoritmi}
Il concetto di relazione è più generale di quello di funzione, in quanto una relazione può associare più elementi del
codominio (es. soluzioni) ad un unico elemento del dominio (es. un'istanza). Una funzione, invece, può associare al massimo
un elemento del codominio ad ogni elemento del dominio.

Un problema computazionale definisce una relazione tra l'insieme delle istanze e quello delle soluzioni in quanto una istanza
può avere più soluzioni valide. Un algoritmo, invece, definisce una funzione tra i due insiemi. A partire da una istanza,
restituirà soltanto una delle possibili soluzioni per tale istanza. È possibile, inoltre, che due algoritmi diversi risolvono
lo stesso problema computazionale, ma producano soluzioni diverse, entrambe valide, a partire dalla stessa istanza.

\subsubsection*{Analisi della correttezza di un algoritmo}
Un algoritmo \(\mathcal{A}\) risolve un problema computazionale \(\pi\) se e solo se, detta \(A\) la funzione di trasferimento
tra ingresso e uscita definita da \(\mathcal{A}\), valgono le seguenti condizioni:
\[1. \quad \mathcal{I}_\text{istanze} \subseteq I_\text{input} \qquad\quad 2. \quad \mathcal{S}_\text{soluzioni} \subseteq O_\text{output} \qquad\quad 3. \quad \forall i \in \mathcal{I} \text{ vale } (i, A(i)) \in \pi\]
\[\text{con} \qquad \pi \subseteq \mathcal{I}_\text{istanze} \times \mathcal{S}_\text{soluzioni} \qquad \text{e} \qquad A: I_\text{input} \to O_\text{output}\]
L'analisi della correttezza consiste in particolare nel verificare la terza condizione.

\subsubsection*{Analisi della complessità di un algoritmo}
L'analisi della complessità computazionale di un algoritmo consiste nel valutare il tempo necessario all'algoritmo per
raggiungere la soluzione a partire da una data istanza generica.

La misurazione del tempo effettivo impiegato dipende molto dall'architettura della macchina su cui avviene la misurazione.
Il valore ottenuto sarebbe troppo variabile e non permetterebbe di stimare in maniera corretta e generale la complessità
di un algoritmo. Per questo motivo si simula l'esecuzione su random access machines con instruction set semplici e tempo
di esecuzione costante per ogni istruzione. Proprio il tempo di esecuzione di una istruzione è usato come unità di misura
del tempo di esecuzione. Questa semplificazione permette di stimare la complessità di un algoritmo in maniera più
generale, misurando banalmente il numero di istruzioni eseguite, e permette inoltre di confrontare tra loro algoritmi diversi
prima di averli implementati su una macchina reale.

Si effettua l'analisi al caso peggiore (worst-case) definendo come tempo impiegato da un algoritmo \(\mathcal{A}\) per
risolvere un'istanza \(i\) di taglia \(n\) come:
\[T_\mathcal{A}(n) \; := \; \sup \;\; T_\mathcal{A}(x) \qquad \text{con} \; \left\{ x \; : \; |x| = n \;\; \wedge \;\; x \in \Sigma^n \right\}\]
L'analisi al caso migliore è poco utile, mentre l'analisi al caso medio è notevolmente più complessa da effettuare, in quanto
richiede di conoscere a priori la distribuzione di probabilità delle varie istanze per una specifica taglia \(n\). In alcuni
casi quest'ultima viene preferita all'analisi al caso peggiore, se il caso peggiore è molto raro e mediamente (per la maggior
parte dei casi) la complessità risulta inferiore, ad esempio per il quicksort o per l'algoritmo del simplesso.

\newpage

\subsection{Dimostrazione per induzione}
\subsubsection*{Assiomi di Peano dei numeri naturali}
\begin{enumerate}
	\item esiste un numero naturale \(0\)
	\item ogni numero naturale \(n\) ha un successore \(s(n)\)
	\item numeri diversi hanno successori diversi: \(n \neq m \implies s(n) \neq s(m)\)
	\item \(0\) non è il successore di nessun numero naturale: \(s(n) \neq 0 \;\; \forall n\)
	\item forma debole del principio di induzione:
	\[\left[ \raisebox{3pt}{ \(\underbrace{P(0)}_{\text{caso base}} \wedge \;\; \underbrace{\forall n \; (P(n) \implies P(n + 1))}_{\text{passo induttivo}} \)} \;\; \right] \implies P(n) \;\; \forall n\]
\end{enumerate}

\vspace{5pt}

\noindent
Il quinto assioma può essere sostituito con il principio del minimo intero che sostiene che ogni insieme ordinabile non
vuoto ha un elemento minimo. Il principio di induzione e quello del minimo intero sono equivalenti, siccome il primo
implica il secondo e viceversa.

\subsubsection*{Forma forte del principio di induzione}
Esiste una formulazione più forte del principio di induzione, equivalente alla forma debole enunciata sopra, ma che facilita
le dimostrazioni basate su questo principio.
\[ \left[ \raisebox{3pt}{ \(\underbrace{P(0)}_{\text{caso base}} \wedge \;\; \underbrace{\forall n \; ((P(0) \wedge P(1) \wedge \ldots \wedge P(n)) \implies P(n + 1))}_{\text{passo induttivo}} \)} \;\; \right] \implies P(n) \;\; \forall n\]
Tale forma forte è ad esempio usata per dimostrare la correttezza del quicksort e negli ordinali transfiniti.

\subsubsection*{Rafforzamento dell'ipotesi induttiva}
Alcune ipotesi induttive più ``permissive'' potrebbero non essere dimostrabili tramite il principio di induzione, anche se
sono vere. Si riescono invece a dimostrare ipotesi più forti e ambiziose che implicano quelle più deboli. Ad esempio
potrebbe non essere dimostrabile che \(T(n) \leq n\), ma si riesce a dimostrare che \(T(n) \leq n - 1\) che implica
\(T(n) \leq n\).

Questo fenomeno paradossale viene chiamato rafforzamento dell'ipotesi induttiva. Si spiega con il fatto che l'ipotesi
induttiva compare sia nelle premesse che nelle conseguenze e, un suo rafforzamento, potrebbe sia rafforzare le conseguenze
che rafforzare le premesse a sua volta. Alcune volte il rafforzamento delle premesse è maggiore rispetto a quello delle
conseguenze e ciò permette di dimostrare queste ultime più facilmente.

\subsubsection*{Mappatura sui numeri naturali}
Il principio di induzione è definito solo sui i numeri naturali e può essere applicato solo su proprietà che hanno come
argomenti numeri naturali. Per dimostrare la correttezza di un algoritmo su un certo insieme di input (che non necessariamente
coincide con l'insieme dei numeri naturali) è necessario effettuare una opportuna mappatura tra gli elementi di tale insieme
e i numeri naturali. Spesso si clusterizza l'insieme degli input in base alla loro taglia e si mappa ogni cluster ad un numero
naturale che corrisponde alla taglia degli input in quel cluster.

Ad esempio, si vorrebbe dimostrare una certa proprietà che vale per gli alberi binari completi, ma gli alberi binari completi
non sono mappabili 1:1 con i numeri naturali. Si clusterizzano, quindi, in base alla loro taglia (ad esempio in base al numero
di nodi interni o alla sua altezza) e si mappa ogni cluster ad un numero naturale equivalente proprio al numero di nodi interni
o alla sua altezza. Si applica quindi il principio di induzione dimostrando che la proprietà vale per ogni cluster.

\subsubsection*{Limiti dell'induzione}
Il principio di induzione si utilizza principalmente per dimostrare ipotesi formulate a priori. Non aiuta a formulare nuove
ipotesi, che vanno pensate con l'aiuto di altre proprietà o intuizioni.

\subsubsection*{Induzione parametrica}
L'induzione parametrica o induzione empirica è una tecnica che permette di individuare i parametri di una funzione generale
o di un modello generale che si ritiene soddisfare al meglio i dati. Può essere usata, ad esempio, nell'addestramento di
reti neurali. In pratica consiste nell'applicare il principio di induzione imponendo che la formula valga per il caso base
e per il passo induttivo.

Ad esempio di vuole calcolare la formula analitica della seguente sommatoria:
\[s(n) = \sum_{k=1}^{n} k^2\]
Si notano le seguenti proprietà:
\begin{itemize}
	\item \(s(n) \leq n^3\) se tutti i valori della sommatoria fossero \(n^2\)
	\item \(s(n) \geq n^3 / 8\) siccome almeno \(n / 2\) valori della sommatoria sono maggiori o uguali a \((n / 2)^2\)
	\item \(s(n) \approx n^3 / 3 - 1/3\) se si tratta la sommatoria come un'integrale finito
\end{itemize}
per cui si ipotizza che la formula analitica sia un'equazione di terzo grado:
\[s(n) = a n^3 + b n^2 + c n + d\]

\begin{itemize}
	\item \textbf{caso base} \\
	si impone che valga il caso base \(s(0) = 0\) da cui si ottiene \(d = 0\)
	\item \textbf{passo induttivo} \\
	si impone che valga il passo induttivo \(s(n + 1) = s(n) + (n + 1)^2\):
\begin{align*}
	a (n + 1)^3 + b (n + 1)^2 + c (n + 1) + d & = a n^3 + b n^2 + c n + d + (n + 1)^2 \\
	a n^3 + (3 a + b) n^2 + (3 a + 2 b + c) n + (a + b + c + d) & = a n^3 + b n^2 + c n + d + n^2 + 2 n + 1 \\
	3a n^2 + (3 a + 2 b) n + (a + b + c) & = n^2 + 2 n + 1
\end{align*}
da cui si ottengono le tre equazioni che determinano gli altri tre parametri \(a\), \(b\) e \(c\):
\[3a = 1 \implies a = 1/3 \qquad\quad 3a + 2b = 2 \implies b = 1/2 \qquad\quad a + b + c = 1 \implies c = 1/6\]
\end{itemize}

\noindent
Si osserva che è possibile trovare gli stessi quattro parametri imponendo che la formula valga per quattro valori
particolari di \(n\) (ad esempio \(n = 0, 1, 2, 3\)) e risolvendo il sistema di quattro equazioni in quattro incognite.
Tale procedimento, però, senza aver osservato all'inizio che la formula è un'equazione di terzo grado, non permette di
affermare che i parametri trovati valgano per ogni \(n\). L'equazione, ad esempio, potrebbe essere di grado superiore
per cui i quattro punti scelti non sarebbero stati sufficienti.

\subsection{Formule matematiche utili}
\vspace{5pt}
\[\begin{array}{c l}
	\displaystyle a \, ^{\log_b c} = c \, ^{\log_b a} & \quad \text{esponenziale di un logaritmo} \\[10pt]
	\displaystyle \sum_{k=0}^{n-1} r^k = \frac{r^n - 1}{r - 1} & \quad \text{serie geometrica di ragione \(r\)} \\[15pt]
	\begin{array}{c}
		\overbrace{z = a + ib}^{\text{cartesiana}} \;\; \leftrightarrow \;\; \overbrace{z = \rho (\cos \varphi + i \sin \varphi)}^{\text{trigonometrica o polare}} \;\; \leftrightarrow \;\; \overbrace{\; z = \rho \cdot e^{i \varphi} \;}^{\substack{\text{esponenziale} \\ \text{o di Eulero}}} \\[5pt]
		\begin{cases}
			\rho = \sqrt{a^2 + b^2} \\
			\varphi = \arctan(b / a)
		\end{cases} \longleftrightarrow \quad \begin{cases}
			a = \rho \cos(\varphi) \\
			b = \rho \sin(\varphi)
		\end{cases}
	\end{array} & \quad \begin{array}{c}
		\text{rappresentazioni di un numero} \\ \text{complesso con passaggio tra} \\ \text{le varie rappresentazioni}
	\end{array} \\[40pt]
	\displaystyle\sum_{k=0}^{n-1} \binom{n}{k} a^k b^{n-k} = (a + b)^n \qquad \sum_{k=0}^{n} \binom{n}{k} = 2^n & \quad \begin{array}{c}
		\text{binomio di Newton con} \\ \text{caso particolare per } a = b = 1 \\ \text{(sottosequenze di una stringa)}
	\end{array} \\[30pt]
\end{array}\]
